\documentclass[12pt,a4paper]{report}
\usepackage[utf8]{inputenc}
\usepackage[T1]{fontenc}
\usepackage[french]{babel}
\usepackage{geometry}
\usepackage{xcolor}
\usepackage{tabularx}
\usepackage{booktabs}
\usepackage{longtable}
\usepackage{enumitem}
\usepackage{fancyhdr}
\usepackage{titlesec}
\usepackage{hyperref}
\usepackage{array}
\usepackage{multirow}
\usepackage{graphicx}
\usepackage{tcolorbox}
\usepackage{mdframed}

\geometry{top=2.5cm, bottom=2.5cm, left=3cm, right=2.5cm}

%----------- Couleurs -----------
\definecolor{primary}{RGB}{30,41,59}
\definecolor{accent}{RGB}{59,130,246}
\definecolor{success}{RGB}{16,185,129}
\definecolor{warning}{RGB}{245,158,11}
\definecolor{danger}{RGB}{239,68,68}
\definecolor{lightgray}{RGB}{248,250,252}
\definecolor{bordergray}{RGB}{226,232,240}

%----------- Hyperlinks -----------
\hypersetup{
    colorlinks=true,
    linkcolor=accent,
    urlcolor=accent,
    pdftitle={Document de Spécification - GRM},
    pdfauthor={Équipe GRM}
}

%----------- En-tête / Pied de page -----------
\pagestyle{fancy}
\fancyhf{}
\fancyhead[L]{\small\color{primary}\textbf{GRM} — Document de Spécification}
\fancyhead[R]{\small\color{primary}\leftmark}
\fancyfoot[L]{\small\color{gray}Confidentiel — Usage interne}
\fancyfoot[C]{\thepage}
\fancyfoot[R]{\small\color{gray}\today}
\renewcommand{\headrulewidth}{0.5pt}
\renewcommand{\footrulewidth}{0.3pt}

%----------- Titres de chapitres -----------
\titleformat{\chapter}[block]
  {\normalfont\Large\bfseries\color{primary}}
  {\thechapter.}{0.8em}{}[\vspace{-4pt}\color{accent}\rule{\textwidth}{2pt}\vspace{6pt}]
\titlespacing*{\chapter}{0pt}{0pt}{16pt}

\titleformat{\section}{\normalfont\large\bfseries\color{primary}}{\thesection}{0.8em}{}
\titleformat{\subsection}{\normalfont\normalsize\bfseries\color{primary}}{\thesubsection}{0.8em}{}

%----------- Boîtes personnalisées -----------
\tcbset{
  exigbox/.style={
    colback=lightgray,
    colframe=bordergray,
    fonttitle=\bfseries\small,
    boxrule=0.5pt,
    arc=4pt,
    left=6pt, right=6pt, top=4pt, bottom=4pt
  }
}

\newcolumntype{L}[1]{>{\raggedright\arraybackslash}p{#1}}
\newcolumntype{C}[1]{>{\centering\arraybackslash}p{#1}}

%==========================================================
\begin{document}
%==========================================================

%----------- PAGE DE GARDE -----------
\begin{titlepage}
\pagecolor{primary}
\color{white}
\vspace*{2cm}
\begin{center}
    {\fontsize{38}{46}\selectfont\bfseries GRM}\\[0.4cm]
    {\Large\bfseries Gestion des Ressources Matérielles}\\[2.5cm]

    \colorbox{accent}{\parbox{0.85\textwidth}{\centering\vspace{0.4cm}
    {\huge\bfseries Document de Spécification}\\[0.2cm]
    {\large Analyse des Besoins \& Cahier des Charges Fonctionnel}
    \vspace{0.4cm}}}

    \vspace{2.5cm}

    \begin{tabular}{ll}
        \textbf{Projet} & Gestion des Ressources Matérielles \\[4pt]
        \textbf{Établissement} & Ynov Campus \\[4pt]
        \textbf{Filière} & M1 Informatique — Architecture Logicielle \\[4pt]
        \textbf{Responsable pédagogique} & Pr. Driss \\[4pt]
        \textbf{Version} & 1.0 \\[4pt]
        \textbf{Date} & \today \\[4pt]
        \textbf{Statut} & \textcolor{success}{Approuvé} \\
    \end{tabular}

    \vspace{3cm}
    {\small\color{gray} Ce document est confidentiel et destiné à un usage interne uniquement.}
\end{center}
\end{titlepage}
\pagecolor{white}\color{black}

%----------- HISTORIQUE DES RÉVISIONS -----------
\chapter*{Historique des Révisions}
\addcontentsline{toc}{chapter}{Historique des Révisions}
\begin{tabularx}{\textwidth}{|C{1.8cm}|C{2.5cm}|L{4.5cm}|L{4.5cm}|}
\hline
\rowcolor{primary}\textcolor{white}{\textbf{Version}} & \textcolor{white}{\textbf{Date}} & \textcolor{white}{\textbf{Description}} & \textcolor{white}{\textbf{Auteur}} \\
\hline
0.1 & 01/04/2025 & Rédaction initiale & Équipe GRM \\
\hline
0.5 & 15/04/2025 & Ajout des cas d'utilisation & Équipe GRM \\
\hline
0.9 & 30/04/2025 & Révision après réunion client & Équipe GRM \\
\hline
1.0 & \today & Version finale approuvée & Équipe GRM \\
\hline
\end{tabularx}

%----------- TABLE DES MATIÈRES -----------
\tableofcontents
\newpage

%==========================================================
\chapter{Introduction Générale}
%==========================================================

\section{Contexte du Projet}

La faculté gère un parc informatique conséquent réparti entre plusieurs départements académiques. Chaque année, l'acquisition de nouvelles ressources matérielles — ordinateurs, imprimantes — fait l'objet de procédures administratives complexes impliquant de nombreux acteurs : enseignants, chefs de département, responsable des ressources, fournisseurs et techniciens de maintenance.

Actuellement, ces processus sont gérés manuellement à travers des échanges de courriels, des tableurs et des documents papier. Cette approche engendre des problèmes récurrents :

\begin{itemize}[leftmargin=1.5cm]
    \item Perte d'informations lors des transmissions entre acteurs
    \item Absence de traçabilité des ressources et de leur historique d'affectation
    \item Délais importants dans le traitement des appels d'offre
    \item Difficultés à suivre l'état des pannes et des interventions de maintenance
    \item Impossibilité de centraliser et de consulter l'inventaire en temps réel
\end{itemize}

\section{Objectifs du Projet}

Le projet \textbf{GRM (Gestion des Ressources Matérielles)} vise à mettre en place un système d'information intégré permettant de :

\begin{enumerate}[leftmargin=1.5cm]
    \item \textbf{Centraliser} la gestion de toutes les ressources matérielles de la faculté
    \item \textbf{Automatiser} le processus des appels d'offre, de la soumission à la sélection
    \item \textbf{Tracer} chaque ressource depuis son acquisition jusqu'à sa mise au rebut
    \item \textbf{Faciliter} la communication entre tous les acteurs du système
    \item \textbf{Sécuriser} les accès aux données selon le rôle de chaque utilisateur
    \item \textbf{Optimiser} la gestion de la maintenance et des pannes
\end{enumerate}

\section{Périmètre du Système}

Le système GRM couvre les domaines fonctionnels suivants :

\begin{tcolorbox}[exigbox, title=Domaines couverts par le système]
\begin{enumerate}
    \item \textbf{Module Appels d'Offre} : création, publication, réception et sélection des offres fournisseurs
    \item \textbf{Module Fournisseurs} : inscription, gestion des profils, liste noire
    \item \textbf{Module Ressources} : inventaire, affectation, suivi du cycle de vie
    \item \textbf{Module Maintenance} : signalement de pannes, constats techniques, retour fournisseur
    \item \textbf{Module Authentification} : gestion des comptes et des accès par rôle
    \item \textbf{Module Départements} : gestion des unités organisationnelles et de leurs membres
\end{enumerate}
\end{tcolorbox}

\section{Public Cible de ce Document}

Ce document de spécification s'adresse aux parties prenantes suivantes :
\begin{itemize}[leftmargin=1.5cm]
    \item L'équipe de développement pour comprendre et implémenter les exigences
    \item Le responsable pédagogique pour valider la conformité au cahier des charges
    \item Les utilisateurs finaux pour valider que leurs besoins ont été correctement compris
\end{itemize}

\section{Conventions de Notation}

\begin{tabularx}{\textwidth}{|C{2cm}|L{10cm}|}
\hline
\rowcolor{lightgray}\textbf{Notation} & \textbf{Signification} \\
\hline
\textbf{EF-XX} & Exigence Fonctionnelle numérotée \\
\hline
\textbf{ENF-XX} & Exigence Non Fonctionnelle numérotée \\
\hline
\textbf{UC-XX} & Cas d'utilisation (Use Case) numéroté \\
\hline
\textbf{RG-XX} & Règle de gestion numérotée \\
\hline
\textbf{[OBLIGATOIRE]} & L'exigence doit être implémentée dans la version courante \\
\hline
\textbf{[OPTIONNEL]} & L'exigence peut être différée à une version ultérieure \\
\hline
\end{tabularx}

%==========================================================
\chapter{Identification des Acteurs}
%==========================================================

\section{Vue d'Ensemble des Acteurs}

Le système GRM implique cinq catégories d'acteurs distincts, chacun disposant de droits d'accès et de responsabilités spécifiques. Le tableau ci-dessous présente une synthèse de ces acteurs.

\begin{longtable}{|L{3.5cm}|L{5.5cm}|L{5.5cm}|}
\caption{Synthèse des acteurs du système} \\
\hline
\rowcolor{primary}\textcolor{white}{\textbf{Acteur}} & \textcolor{white}{\textbf{Responsabilités principales}} & \textcolor{white}{\textbf{Interactions avec le système}} \\
\hline
\endfirsthead
\hline
\rowcolor{primary}\textcolor{white}{\textbf{Acteur}} & \textcolor{white}{\textbf{Responsabilités principales}} & \textcolor{white}{\textbf{Interactions avec le système}} \\
\hline
\endhead
Responsable des ressources &
Administrateur du parc matériel. Supervise les appels d'offre, sélectionne les fournisseurs, gère l'inventaire, décide des actions de maintenance. &
Créer des AO, sélectionner offres, affecter ressources, blacklister fournisseurs, valider constats de maintenance. \\
\hline
Chef de département &
Représente un département académique. Collecte les besoins des enseignants et les soumet au responsable. Reçoit les affectations. &
Saisir les besoins, suivre les affectations, signaler des pannes, consulter les ressources du département. \\
\hline
Enseignant &
Utilisateur final des ressources matérielles. Exprime ses besoins et signale les dysfonctionnements. &
Exprimer besoins (via chef), signaler pannes, consulter ressources affectées. \\
\hline
Fournisseur (Société) &
Entreprise externe souhaitant répondre aux appels d'offre de la faculté. &
S'inscrire, consulter AO ouverts, soumettre offres avec détail prix/garantie/livraison. \\
\hline
Technicien de maintenance &
Agent technique chargé du diagnostic et de la réparation des ressources défectueuses. &
Consulter les pannes signalées, intervenir sur site, rédiger des constats techniques détaillés. \\
\hline
\end{longtable}

\section{Fiche Détaillée des Acteurs}

\subsection{Responsable des Ressources}

\begin{tcolorbox}[exigbox, title=Responsable des Ressources]
\textbf{Rôle système :} \texttt{responsable} \\
\textbf{Niveau d'accès :} Maximum — accès à l'intégralité des fonctionnalités \\
\textbf{Profil type :} Administrateur informatique ou responsable logistique de la faculté \\[6pt]
\textbf{Besoins métier :}
\begin{itemize}
    \item Disposer d'une vue consolidée de tous les besoins des départements avant de lancer un appel d'offre
    \item Pouvoir comparer objectivement les offres des fournisseurs (prix, délai, garantie)
    \item Maintenir un inventaire précis avec numéros de série et historique des affectations
    \item Suivre l'état de toutes les pannes et décider des actions correctives
    \item Protéger la faculté contre les fournisseurs défaillants via la liste noire
\end{itemize}
\end{tcolorbox}

\subsection{Chef de Département}

\begin{tcolorbox}[exigbox, title=Chef de Département]
\textbf{Rôle système :} \texttt{chef\_departement} \\
\textbf{Niveau d'accès :} Intermédiaire — limité aux ressources de son département \\
\textbf{Profil type :} Professeur responsable d'une unité pédagogique \\[6pt]
\textbf{Besoins métier :}
\begin{itemize}
    \item Recueillir les besoins de son équipe et les soumettre de manière structurée
    \item Consulter la liste des ressources affectées à son département
    \item Participer aux réunions de concertation avant validation des besoins
    \item Être informé des affectations prévues suite à un appel d'offre
\end{itemize}
\end{tcolorbox}

\subsection{Enseignant}

\begin{tcolorbox}[exigbox, title=Enseignant]
\textbf{Rôle système :} \texttt{enseignant} \\
\textbf{Niveau d'accès :} Limité — consultation et signalement uniquement \\
\textbf{Profil type :} Membre du corps enseignant d'un département \\[6pt]
\textbf{Besoins métier :}
\begin{itemize}
    \item Signaler rapidement toute panne affectant son travail
    \item Consulter les ressources qui lui sont affectées
    \item Suivre le statut de ses signalements de pannes
\end{itemize}
\end{tcolorbox}

\subsection{Fournisseur}

\begin{tcolorbox}[exigbox, title=Fournisseur]
\textbf{Rôle système :} \texttt{fournisseur} \\
\textbf{Niveau d'accès :} Externe — consultation AO ouverts et soumission d'offres \\
\textbf{Profil type :} Société commerciale spécialisée en matériel informatique \\[6pt]
\textbf{Besoins métier :}
\begin{itemize}
    \item Accéder facilement aux appels d'offre en cours sans démarches administratives lourdes
    \item Soumettre des propositions détaillées avec tarification article par article
    \item Être informé rapidement de l'acceptation ou du rejet de son offre avec motif
\end{itemize}
\end{tcolorbox}

\subsection{Technicien de Maintenance}

\begin{tcolorbox}[exigbox, title=Technicien de Maintenance]
\textbf{Rôle système :} \texttt{technicien} \\
\textbf{Niveau d'accès :} Limité — consultation des pannes et rédaction de constats \\
\textbf{Profil type :} Agent de support technique de la faculté \\[6pt]
\textbf{Besoins métier :}
\begin{itemize}
    \item Consulter la liste des pannes ouvertes pour planifier ses interventions
    \item Rédiger un constat structuré après chaque intervention
    \item Indiquer si la ressource peut être réparée sur place ou doit être retournée
\end{itemize}
\end{tcolorbox}

%==========================================================
\chapter{Processus Métier}
%==========================================================

\section{Vue Générale des Processus}

Le système GRM s'articule autour de trois processus métier principaux, chacun impliquant plusieurs acteurs et pouvant être modélisé en notation BPMN.

\section{Processus 1 : Gestion des Appels d'Offre}

\subsection{Description Générale}

Ce processus couvre le cycle complet depuis l'identification des besoins jusqu'à l'attribution du marché à un fournisseur.

\subsection{Déroulement Détaillé}

\begin{longtable}{|C{0.8cm}|L{3cm}|L{6cm}|L{4cm}|}
\caption{Étapes du processus Appel d'Offre} \\
\hline
\rowcolor{primary}\textcolor{white}{\textbf{\#}} & \textcolor{white}{\textbf{Acteur}} & \textcolor{white}{\textbf{Action}} & \textcolor{white}{\textbf{Résultat}} \\
\hline
\endfirsthead
\hline
\textbf{\#} & \textbf{Acteur} & \textbf{Action} & \textbf{Résultat} \\
\hline
\endhead
1 & Chef de département & Interroger les enseignants sur leurs besoins en matériel & Liste des besoins bruts \\
\hline
2 & Enseignants & Exprimer leurs besoins (type, quantité, spécifications) & Besoins individuels collectés \\
\hline
3 & Chef de département & Organiser une réunion de concertation pour valider et ajuster les besoins & Besoins validés et consolidés \\
\hline
4 & Chef de département & Soumettre les besoins consolidés au responsable via le système & Demande formalisée \\
\hline
5 & Responsable & Rassembler les besoins de tous les départements & Besoins globaux de la faculté \\
\hline
6 & Responsable & Créer l'appel d'offre avec dates de début et de fin & AO publié au statut «ouvert» \\
\hline
7 & Fournisseurs & Consulter les appels d'offre ouverts & Connaissance des opportunités \\
\hline
8 & Fournisseurs & Soumettre une offre (prix, garantie, délai, détail) & Offre enregistrée \\
\hline
9 & Responsable & Vérifier et éliminer les fournisseurs de la liste noire & Offres non conformes filtrées \\
\hline
10 & Responsable & Sélectionner l'offre la moins disante parmi les conformes & Fournisseur retenu \\
\hline
11 & Système & Envoyer notification d'acceptation au fournisseur retenu & Fournisseur informé \\
\hline
12 & Système & Envoyer notifications de rejet aux autres fournisseurs & Candidats non retenus informés \\
\hline
13 & Responsable & Fermer l'appel d'offre et attendre la livraison & AO au statut «attribué» \\
\hline
\end{longtable}

\subsection{Points de Décision}

\begin{itemize}[leftmargin=1.5cm]
    \item \textbf{Réunion de concertation} : Le chef peut modifier les besoins après consultation. Les modifications sont possibles jusqu'à soumission au responsable.
    \item \textbf{Filtrage des offres} : Le responsable peut rejeter une offre individuellement avant sélection finale.
    \item \textbf{Sélection} : Si aucune offre n'est satisfaisante, le responsable peut laisser l'AO ouvert ou le recréer.
\end{itemize}

\section{Processus 2 : Réception et Gestion de l'Inventaire}

\subsection{Description Générale}

Après attribution du marché, ce processus gère la réception physique des ressources, leur enregistrement dans l'inventaire et leur affectation aux départements.

\subsection{Déroulement Détaillé}

\begin{longtable}{|C{0.8cm}|L{3cm}|L{6cm}|L{4cm}|}
\caption{Étapes du processus Réception et Inventaire} \\
\hline
\rowcolor{primary}\textcolor{white}{\textbf{\#}} & \textcolor{white}{\textbf{Acteur}} & \textcolor{white}{\textbf{Action}} & \textcolor{white}{\textbf{Résultat}} \\
\hline
\endfirsthead
\hline
\textbf{\#} & \textbf{Acteur} & \textbf{Action} & \textbf{Résultat} \\
\hline
\endhead
1 & Fournisseur & Livrer les ressources matérielles à la faculté & Matériel reçu \\
\hline
2 & Responsable & Vérifier la conformité des ressources livrées & Réception validée ou litige \\
\hline
3 & Responsable & Si nouveau fournisseur : saisir lieu, adresse, site, gérant & Fiche fournisseur complète \\
\hline
4 & Responsable & Attribuer un numéro d'inventaire (code à barres) à chaque ressource & Ressource enregistrée \\
\hline
5 & Responsable & Enregistrer les caractéristiques techniques de chaque ressource & Fiche technique créée \\
\hline
6 & Responsable & Enregistrer la date de fin de garantie de chaque ressource & Garantie tracée \\
\hline
7 & Responsable & Affecter chaque ressource à un département (ou une personne précise) & Affectation enregistrée \\
\hline
8 & Système & Mettre à jour le statut de la ressource («affectée») & Inventaire à jour \\
\hline
9 & Système & Générer et envoyer la liste des affectations au chef de département & Chef informé \\
\hline
\end{longtable}

\section{Processus 3 : Maintenance et Traitement des Pannes}

\subsection{Description Générale}

Ce processus gère le cycle de vie d'une panne, depuis son signalement par un enseignant jusqu'à sa résolution ou le retour de la ressource au fournisseur.

\subsection{Déroulement Détaillé}

\begin{longtable}{|C{0.8cm}|L{3cm}|L{6cm}|L{4cm}|}
\caption{Étapes du processus Maintenance} \\
\hline
\rowcolor{primary}\textcolor{white}{\textbf{\#}} & \textcolor{white}{\textbf{Acteur}} & \textcolor{white}{\textbf{Action}} & \textcolor{white}{\textbf{Résultat}} \\
\hline
\endfirsthead
\hline
\textbf{\#} & \textbf{Acteur} & \textbf{Action} & \textbf{Résultat} \\
\hline
\endhead
1 & Enseignant / Chef & Signaler une panne sur une ressource de son département & Ticket de panne créé \\
\hline
2 & Système & Changer le statut de la ressource à «maintenance» & Ressource marquée hors service \\
\hline
3 & Technicien & Consulter les pannes ouvertes et planifier son intervention & Planning de maintenance \\
\hline
4 & Technicien & Intervenir physiquement auprès du département concerné & Diagnostic effectué \\
\hline
5a & Technicien & Si panne simple : réparer la ressource sur place & Panne résolue \\
\hline
5b & Technicien & Si panne sévère : rédiger un constat (explication, date, fréquence, type, ordre logiciel ou matériel) & Constat transmis \\
\hline
6 & Système & Transmettre le constat au responsable des ressources & Responsable alerté \\
\hline
7a & Responsable & Si réparable et garantie valide : renvoyer au fournisseur pour réparation & Retour fournisseur engagé \\
\hline
7b & Responsable & Si garantie expirée ou non réparable : décision de remplacement & Nouveau besoin créé \\
\hline
8 & Responsable & Résoudre le ticket de panne et remettre la ressource en service & Ressource de nouveau disponible \\
\hline
\end{longtable}

\subsection{Règle Spécifique — Imprimantes}

\begin{tcolorbox}[exigbox, title=Règle de Gestion RG-01]
Les pannes signalées sur des imprimantes sont \textbf{systématiquement} considérées comme d'ordre matériel. Le champ «type d'ordre» pour une imprimante ne peut prendre que la valeur «matériel». Cette règle est appliquée côté serveur lors de la création du constat.
\end{tcolorbox}

%==========================================================
\chapter{Exigences Fonctionnelles}
%==========================================================

\section{Module Authentification et Sécurité}

\begin{longtable}{|C{1.8cm}|L{8.5cm}|C{2cm}|C{2cm}|}
\caption{Exigences — Module Authentification} \\
\hline
\rowcolor{primary}\textcolor{white}{\textbf{ID}} & \textcolor{white}{\textbf{Description}} & \textcolor{white}{\textbf{Priorité}} & \textcolor{white}{\textbf{Statut}} \\
\hline
\endfirsthead
\hline
\textbf{ID} & \textbf{Description} & \textbf{Priorité} & \textbf{Statut} \\
\hline
\endhead
EF-AUTH-01 & Le système doit permettre à tout utilisateur de s'authentifier via une adresse e-mail et un mot de passe. & Haute & Impl. \\
\hline
EF-AUTH-02 & Le système doit retourner un jeton JWT signé à l'authentification réussie, valide pendant 24 heures. & Haute & Impl. \\
\hline
EF-AUTH-03 & Le système doit rejeter toute requête sur une ressource protégée si aucun jeton valide n'est fourni (HTTP 401). & Haute & Impl. \\
\hline
EF-AUTH-04 & Le système doit rejeter tout accès à une fonctionnalité non autorisée pour le rôle de l'utilisateur connecté (HTTP 403). & Haute & Impl. \\
\hline
EF-AUTH-05 & Un fournisseur doit pouvoir s'inscrire en fournissant le nom de sa société sans nécessiter de validation administrative préalable. & Haute & Impl. \\
\hline
EF-AUTH-06 & L'utilisateur doit pouvoir consulter son profil (nom, rôle, département). & Moyenne & Impl. \\
\hline
\end{longtable}

\section{Module Appels d'Offre}

\begin{longtable}{|C{1.8cm}|L{8.5cm}|C{2cm}|C{2cm}|}
\caption{Exigences — Module Appels d'Offre} \\
\hline
\rowcolor{primary}\textcolor{white}{\textbf{ID}} & \textcolor{white}{\textbf{Description}} & \textcolor{white}{\textbf{Priorité}} & \textcolor{white}{\textbf{Statut}} \\
\hline
\endfirsthead
\hline
\textbf{ID} & \textbf{Description} & \textbf{Priorité} & \textbf{Statut} \\
\hline
\endhead
EF-AO-01 & Le responsable et le chef de département doivent pouvoir créer un appel d'offre avec titre, description, dates de début et de fin. & Haute & Impl. \\
\hline
EF-AO-02 & Un appel d'offre doit pouvoir contenir plusieurs besoins en ressources, chacun spécifiant le type (ordinateur/imprimante), la quantité et les spécifications techniques. & Haute & Impl. \\
\hline
EF-AO-03 & Pour un ordinateur, les spécifications saisissables sont : marque, CPU, RAM, disque dur, écran. & Haute & Impl. \\
\hline
EF-AO-04 & Pour une imprimante, les spécifications saisissables sont : marque, vitesse d'impression, résolution. & Haute & Impl. \\
\hline
EF-AO-05 & La liste de tous les appels d'offre doit être consultable par tous les utilisateurs authentifiés. & Haute & Impl. \\
\hline
EF-AO-06 & Les fournisseurs authentifiés doivent pouvoir consulter uniquement les appels d'offre au statut «ouvert». & Haute & Impl. \\
\hline
EF-AO-07 & Le responsable doit pouvoir fermer un appel d'offre, passant son statut de «ouvert» à «fermé». & Haute & Impl. \\
\hline
EF-AO-08 & Un appel d'offre passe automatiquement au statut «attribué» lorsqu'une offre est acceptée. & Haute & Impl. \\
\hline
EF-AO-09 & Le responsable doit pouvoir modifier un appel d'offre tant qu'il est au statut «brouillon» ou «ouvert». & Moyenne & Impl. \\
\hline
\end{longtable}

\section{Module Gestion des Offres Fournisseurs}

\begin{longtable}{|C{1.8cm}|L{8.5cm}|C{2cm}|C{2cm}|}
\caption{Exigences — Module Offres} \\
\hline
\rowcolor{primary}\textcolor{white}{\textbf{ID}} & \textcolor{white}{\textbf{Description}} & \textcolor{white}{\textbf{Priorité}} & \textcolor{white}{\textbf{Statut}} \\
\hline
\endfirsthead
\hline
\textbf{ID} & \textbf{Description} & \textbf{Priorité} & \textbf{Statut} \\
\hline
\endhead
EF-OFF-01 & Un fournisseur non blacklisté doit pouvoir soumettre une offre sur un appel d'offre ouvert. & Haute & Impl. \\
\hline
EF-OFF-02 & Une offre doit contenir obligatoirement : la date de livraison prévue, la durée de garantie en mois, le prix total. & Haute & Impl. \\
\hline
EF-OFF-03 & Une offre doit permettre de détailler le prix unitaire, la marque et la quantité pour chaque type de ressource proposé. & Haute & Impl. \\
\hline
EF-OFF-04 & Le système doit refuser toute soumission d'offre par un fournisseur blacklisté, même si ce dernier dispose d'un compte valide. & Haute & Impl. \\
\hline
EF-OFF-05 & Le responsable doit pouvoir visualiser toutes les offres d'un appel d'offre, triées par prix total croissant. & Haute & Impl. \\
\hline
EF-OFF-06 & Le responsable doit pouvoir accepter une offre ; le système rejette automatiquement toutes les autres offres du même appel. & Haute & Impl. \\
\hline
EF-OFF-07 & Le responsable doit pouvoir rejeter une offre individuellement en saisissant un motif de rejet. & Haute & Impl. \\
\hline
EF-OFF-08 & Le fournisseur dont l'offre est rejetée doit recevoir une notification incluant le motif du rejet. & Moyenne & Impl. \\
\hline
EF-OFF-09 & Le fournisseur dont l'offre est acceptée doit recevoir une notification d'acceptation. & Moyenne & Impl. \\
\hline
\end{longtable}

\section{Module Gestion des Fournisseurs}

\begin{longtable}{|C{1.8cm}|L{8.5cm}|C{2cm}|C{2cm}|}
\caption{Exigences — Module Fournisseurs} \\
\hline
\rowcolor{primary}\textcolor{white}{\textbf{ID}} & \textcolor{white}{\textbf{Description}} & \textcolor{white}{\textbf{Priorité}} & \textcolor{white}{\textbf{Statut}} \\
\hline
\endfirsthead
\hline
\textbf{ID} & \textbf{Description} & \textbf{Priorité} & \textbf{Statut} \\
\hline
\endhead
EF-FRN-01 & L'inscription d'un fournisseur nécessite uniquement : nom de la société, nom du contact, adresse e-mail, mot de passe. & Haute & Impl. \\
\hline
EF-FRN-02 & Le profil fournisseur peut être enrichi avec : ville/lieu, adresse postale, site internet, nom du gérant. & Moyenne & Impl. \\
\hline
EF-FRN-03 & Le responsable doit pouvoir consulter la liste complète de tous les fournisseurs inscrits. & Haute & Impl. \\
\hline
EF-FRN-04 & Le responsable doit pouvoir mettre un fournisseur en liste noire en obligeant la saisie d'un motif. & Haute & Impl. \\
\hline
EF-FRN-05 & Le responsable doit pouvoir retirer un fournisseur de la liste noire. & Haute & Impl. \\
\hline
EF-FRN-06 & Le statut de liste noire d'un fournisseur doit être clairement visible dans l'interface. & Haute & Impl. \\
\hline
EF-FRN-07 & L'historique des offres soumises par un fournisseur doit être consultable par le responsable. & Moyenne & Impl. \\
\hline
\end{longtable}

\section{Module Gestion des Ressources}

\begin{longtable}{|C{1.8cm}|L{8.5cm}|C{2cm}|C{2cm}|}
\caption{Exigences — Module Ressources} \\
\hline
\rowcolor{primary}\textcolor{white}{\textbf{ID}} & \textcolor{white}{\textbf{Description}} & \textcolor{white}{\textbf{Priorité}} & \textcolor{white}{\textbf{Statut}} \\
\hline
\endfirsthead
\hline
\textbf{ID} & \textbf{Description} & \textbf{Priorité} & \textbf{Statut} \\
\hline
\endhead
EF-RES-01 & Le responsable doit pouvoir créer une fiche ressource avec numéro d'inventaire unique, type, marque et date de garantie. & Haute & Impl. \\
\hline
EF-RES-02 & Le numéro d'inventaire (code à barres) doit être unique dans tout le système. & Haute & Impl. \\
\hline
EF-RES-03 & Le responsable doit pouvoir lister toutes les ressources avec filtrage par type et statut. & Haute & Impl. \\
\hline
EF-RES-04 & Le responsable doit pouvoir modifier les informations d'une ressource existante. & Haute & Impl. \\
\hline
EF-RES-05 & Le responsable doit pouvoir supprimer une ressource de l'inventaire. & Haute & Impl. \\
\hline
EF-RES-06 & Le responsable doit pouvoir affecter une ressource à un département entier ou à un enseignant précis. & Haute & Impl. \\
\hline
EF-RES-07 & Une nouvelle affectation doit automatiquement désactiver l'affectation précédente (historique conservé). & Haute & Impl. \\
\hline
EF-RES-08 & Le statut d'une ressource évolue automatiquement : disponible → affectée → maintenance → retournée. & Haute & Impl. \\
\hline
EF-RES-09 & L'historique complet des affectations d'une ressource doit être consultable. & Haute & Impl. \\
\hline
EF-RES-10 & Les spécifications techniques d'une ressource (CPU, RAM, etc.) doivent être consultables sur sa fiche détaillée. & Moyenne & Impl. \\
\hline
\end{longtable}

\section{Module Maintenance et Pannes}

\begin{longtable}{|C{1.8cm}|L{8.5cm}|C{2cm}|C{2cm}|}
\caption{Exigences — Module Maintenance} \\
\hline
\rowcolor{primary}\textcolor{white}{\textbf{ID}} & \textcolor{white}{\textbf{Description}} & \textcolor{white}{\textbf{Priorité}} & \textcolor{white}{\textbf{Statut}} \\
\hline
\endfirsthead
\hline
\textbf{ID} & \textbf{Description} & \textbf{Priorité} & \textbf{Statut} \\
\hline
\endhead
EF-MNT-01 & Un enseignant ou chef de département doit pouvoir signaler une panne en sélectionnant la ressource et en décrivant le dysfonctionnement. & Haute & Impl. \\
\hline
EF-MNT-02 & Le signalement d'une panne doit automatiquement changer le statut de la ressource à «maintenance». & Haute & Impl. \\
\hline
EF-MNT-03 & Le technicien doit pouvoir consulter toutes les pannes ouvertes qui lui sont assignées. & Haute & Impl. \\
\hline
EF-MNT-04 & Le technicien doit pouvoir rédiger un constat de panne contenant : explication, date d'apparition, fréquence (rare/fréquente/permanente), type d'ordre (logiciel/matériel). & Haute & Impl. \\
\hline
EF-MNT-05 & Le technicien doit indiquer si la ressource peut ou non être réparée sur place. & Haute & Impl. \\
\hline
EF-MNT-06 & Le responsable doit recevoir le constat du technicien et pouvoir décider d'un retour fournisseur. & Haute & Impl. \\
\hline
EF-MNT-07 & Si la ressource est retournée au fournisseur, son statut doit passer à «retournée». & Haute & Impl. \\
\hline
EF-MNT-08 & Le responsable doit pouvoir résoudre un ticket de panne, remettant la ressource au statut «disponible». & Haute & Impl. \\
\hline
EF-MNT-09 & La date de fin de garantie de la ressource doit être affichée lors de la consultation du ticket, pour faciliter la décision du responsable. & Haute & Impl. \\
\hline
EF-MNT-10 & Le système doit distinguer visuellement si la garantie est encore valide ou expirée. & Moyenne & Impl. \\
\hline
\end{longtable}

%==========================================================
\chapter{Exigences Non Fonctionnelles}
%==========================================================

\section{Sécurité}

\begin{longtable}{|C{1.8cm}|L{9cm}|C{2cm}|}
\caption{Exigences de sécurité} \\
\hline
\rowcolor{primary}\textcolor{white}{\textbf{ID}} & \textcolor{white}{\textbf{Description}} & \textcolor{white}{\textbf{Priorité}} \\
\hline
\endfirsthead
\hline
\textbf{ID} & \textbf{Description} & \textbf{Priorité} \\
\hline
\endhead
ENF-SEC-01 & Tous les mots de passe doivent être hachés avec l'algorithme bcrypt (facteur de coût minimum : 10) avant stockage. Aucun mot de passe en clair ne doit jamais être persisté. & Critique \\
\hline
ENF-SEC-02 & La communication entre le client et le serveur doit transiter exclusivement via HTTPS en production. & Critique \\
\hline
ENF-SEC-03 & Les tokens JWT doivent être signés avec un secret d'au moins 256 bits, configuré via variable d'environnement (non codé en dur). & Critique \\
\hline
ENF-SEC-04 & Les tokens JWT ont une durée de validité de 24 heures. Passé ce délai, l'utilisateur doit se réauthentifier. & Haute \\
\hline
ENF-SEC-05 & Toutes les requêtes vers l'API doivent vérifier la présence et la validité du token JWT avant traitement. & Critique \\
\hline
ENF-SEC-06 & Le contrôle d'accès par rôle (RBAC) doit être appliqué côté serveur pour chaque endpoint. Une vérification côté client seule ne suffit pas. & Critique \\
\hline
ENF-SEC-07 & Les requêtes SQL doivent utiliser des requêtes paramétrées (prepared statements) pour prévenir les injections SQL. & Critique \\
\hline
ENF-SEC-08 & Les en-têtes CORS doivent être configurés de manière restrictive en production pour n'autoriser que les origines légitimes. & Haute \\
\hline
\end{longtable}

\section{Performance}

\begin{longtable}{|C{1.8cm}|L{9cm}|C{2cm}|}
\caption{Exigences de performance} \\
\hline
\rowcolor{primary}\textcolor{white}{\textbf{ID}} & \textcolor{white}{\textbf{Description}} & \textcolor{white}{\textbf{Priorité}} \\
\hline
\endfirsthead
\hline
\textbf{ID} & \textbf{Description} & \textbf{Priorité} \\
\hline
\endhead
ENF-PERF-01 & Le temps de réponse de l'API doit être inférieur à 500 ms pour 95\% des requêtes en conditions normales d'utilisation (< 50 utilisateurs simultanés). & Haute \\
\hline
ENF-PERF-02 & Le chargement initial de l'application React doit être inférieur à 3 secondes sur une connexion de 10 Mbps. & Haute \\
\hline
ENF-PERF-03 & Les requêtes de liste (ressources, appels d'offre) doivent supporter la pagination pour les ensembles de données > 1000 éléments. & Moyenne \\
\hline
ENF-PERF-04 & Le pool de connexions PostgreSQL doit être correctement dimensionné (min 5, max 20 connexions) pour éviter les saturations. & Haute \\
\hline
\end{longtable}

\section{Disponibilité et Fiabilité}

\begin{longtable}{|C{1.8cm}|L{9cm}|C{2cm}|}
\caption{Exigences de disponibilité} \\
\hline
\rowcolor{primary}\textcolor{white}{\textbf{ID}} & \textcolor{white}{\textbf{Description}} & \textcolor{white}{\textbf{Priorité}} \\
\hline
\endfirsthead
\hline
\textbf{ID} & \textbf{Description} & \textbf{Priorité} \\
\hline
\endhead
ENF-DISPO-01 & Le système doit être disponible 99,5\% du temps pendant les heures ouvrables (8h-20h, du lundi au vendredi). & Haute \\
\hline
ENF-DISPO-02 & Les opérations critiques (acceptation d'offre, affectation de ressource) doivent être gérées dans des transactions SQL atomiques (BEGIN/COMMIT/ROLLBACK). & Critique \\
\hline
ENF-DISPO-03 & En cas d'erreur serveur, l'API doit retourner des messages d'erreur explicites (code HTTP + message JSON) sans exposer de détails techniques internes. & Haute \\
\hline
ENF-DISPO-04 & La base de données doit être sauvegardée quotidiennement avec une rétention d'au moins 7 jours. & Haute \\
\hline
\end{longtable}

\section{Maintenabilité et Évolutivité}

\begin{longtable}{|C{1.8cm}|L{9cm}|C{2cm}|}
\caption{Exigences de maintenabilité} \\
\hline
\rowcolor{primary}\textcolor{white}{\textbf{ID}} & \textcolor{white}{\textbf{Description}} & \textcolor{white}{\textbf{Priorité}} \\
\hline
\endfirsthead
\hline
\textbf{ID} & \textbf{Description} & \textbf{Priorité} \\
\hline
\endhead
ENF-MAINT-01 & L'architecture doit séparer clairement les couches présentation, logique métier et persistance des données. & Haute \\
\hline
ENF-MAINT-02 & Le code source doit être versionné sur une plateforme Git (GitHub/GitLab). & Haute \\
\hline
ENF-MAINT-03 & Toutes les variables de configuration (connexion DB, JWT secret, port) doivent être externalisées dans des fichiers \texttt{.env}, non versionnés. & Haute \\
\hline
ENF-MAINT-04 & Le déploiement doit pouvoir être reproduit intégralement via Docker Compose en une seule commande. & Haute \\
\hline
ENF-MAINT-05 & L'application doit exposer un endpoint de santé (\texttt{GET /api/health}) permettant aux outils de monitoring de vérifier sa disponibilité. & Moyenne \\
\hline
\end{longtable}

%==========================================================
\chapter{Règles de Gestion}
%==========================================================

\begin{longtable}{|C{1.8cm}|L{4cm}|L{8cm}|}
\caption{Règles de gestion métier} \\
\hline
\rowcolor{primary}\textcolor{white}{\textbf{ID}} & \textcolor{white}{\textbf{Domaine}} & \textcolor{white}{\textbf{Règle}} \\
\hline
\endfirsthead
\hline
\textbf{ID} & \textbf{Domaine} & \textbf{Règle} \\
\hline
\endhead
RG-01 & Imprimantes & Une panne d'imprimante est toujours d'ordre matériel. Le champ «ordre logiciel» ne s'applique pas aux imprimantes. \\
\hline
RG-02 & Appel d'offre & Lorsqu'une offre est acceptée, toutes les autres offres du même appel d'offre sont automatiquement rejetées. \\
\hline
RG-03 & Appel d'offre & L'appel d'offre passe au statut «attribué» dès qu'une offre est acceptée. Aucune autre offre ne peut être acceptée. \\
\hline
RG-04 & Fournisseur & Un fournisseur blacklisté ne peut pas soumettre d'offre. Le contrôle s'effectue au moment de la soumission, même si le blacklist est postérieur à l'inscription. \\
\hline
RG-05 & Affectation & Lorsqu'une ressource est affectée à une nouvelle destination, son affectation précédente est automatiquement clôturée (date de retour enregistrée). \\
\hline
RG-06 & Inventaire & Le numéro d'inventaire d'une ressource est unique dans tout le système et ne peut pas être réutilisé. \\
\hline
RG-07 & Maintenance & La signalement d'une panne change automatiquement le statut de la ressource à «maintenance». \\
\hline
RG-08 & Maintenance & La résolution d'une panne remet automatiquement le statut de la ressource à «disponible». \\
\hline
RG-09 & Maintenance & Le retour d'une ressource au fournisseur change son statut à «retournée». Elle ne peut plus être affectée jusqu'à décision du responsable. \\
\hline
RG-10 & Sécurité & Toute action sensible (acceptation offre, blacklist, suppression ressource) est réservée exclusivement au rôle «responsable». \\
\hline
\end{longtable}

%==========================================================
\chapter{Modèle Conceptuel des Données}
%==========================================================

\section{Entités et Attributs}

\subsection{Entité Utilisateur}

\begin{tabularx}{\textwidth}{|L{3cm}|L{3cm}|L{7cm}|}
\hline
\rowcolor{lightgray}\textbf{Attribut} & \textbf{Type} & \textbf{Description} \\
\hline
id & UUID & Identifiant unique généré automatiquement \\
\hline
name & VARCHAR(255) & Nom complet de l'utilisateur \\
\hline
email & VARCHAR(255) & Adresse e-mail, unique dans le système \\
\hline
password & VARCHAR(255) & Mot de passe haché (bcrypt) \\
\hline
role & ENUM & Rôle : responsable, chef\_departement, enseignant, fournisseur, technicien \\
\hline
department\_id & UUID (FK) & Référence au département (nullable pour responsable et technicien) \\
\hline
created\_at & TIMESTAMP & Date de création du compte \\
\hline
\end{tabularx}

\subsection{Entité Appel d'Offre (Tender)}

\begin{tabularx}{\textwidth}{|L{3cm}|L{3cm}|L{7cm}|}
\hline
\rowcolor{lightgray}\textbf{Attribut} & \textbf{Type} & \textbf{Description} \\
\hline
id & UUID & Identifiant unique \\
\hline
title & VARCHAR(255) & Intitulé de l'appel d'offre \\
\hline
description & TEXT & Description générale \\
\hline
start\_date & DATE & Date d'ouverture des offres \\
\hline
end\_date & DATE & Date de clôture des offres \\
\hline
status & ENUM & draft / open / closed / awarded \\
\hline
created\_by & UUID (FK) & Utilisateur créateur \\
\hline
created\_at & TIMESTAMP & Date de création \\
\hline
\end{tabularx}

\subsection{Entité Ressource}

\begin{tabularx}{\textwidth}{|L{3cm}|L{3cm}|L{7cm}|}
\hline
\rowcolor{lightgray}\textbf{Attribut} & \textbf{Type} & \textbf{Description} \\
\hline
id & UUID & Identifiant unique \\
\hline
inventory\_number & VARCHAR(100) & Numéro d'inventaire / code à barres (unique) \\
\hline
resource\_type & ENUM & ordinateur / imprimante \\
\hline
brand & VARCHAR(255) & Marque du matériel \\
\hline
status & ENUM & available / assigned / maintenance / returned \\
\hline
warranty\_end & DATE & Date de fin de garantie \\
\hline
specs & JSONB & Spécifications techniques (CPU, RAM, résolution…) \\
\hline
supplier\_id & UUID (FK) & Fournisseur d'origine \\
\hline
offer\_id & UUID (FK) & Offre d'acquisition \\
\hline
\end{tabularx}

\section{Relations Principales}

\begin{tabularx}{\textwidth}{|L{3.5cm}|C{2cm}|L{3.5cm}|L{5cm}|}
\hline
\rowcolor{lightgray}\textbf{Entité A} & \textbf{Cardinalité} & \textbf{Entité B} & \textbf{Description} \\
\hline
Département & 1..* & Utilisateur & Un département a plusieurs membres \\
\hline
AppelOffre & 1..* & BesoinRessource & Un AO a plusieurs besoins \\
\hline
AppelOffre & 0..* & Offre & Un AO reçoit plusieurs offres \\
\hline
Fournisseur & 1..* & Offre & Un fournisseur soumet plusieurs offres \\
\hline
Offre & 1..* & ArticleOffre & Une offre détaille plusieurs articles \\
\hline
Ressource & 0..* & Affectation & Une ressource a plusieurs affectations (historique) \\
\hline
Ressource & 0..* & Panne & Une ressource peut avoir plusieurs pannes \\
\hline
Panne & 0..1 & ConstatMaintenance & Une panne peut avoir un constat technique \\
\hline
\end{tabularx}

%==========================================================
\chapter{Glossaire}
%==========================================================

\begin{longtable}{|L{3.5cm}|L{10cm}|}
\caption{Glossaire des termes métier} \\
\hline
\rowcolor{primary}\textcolor{white}{\textbf{Terme}} & \textcolor{white}{\textbf{Définition}} \\
\hline
\endfirsthead
\hline
\textbf{Terme} & \textbf{Définition} \\
\hline
\endhead
Appel d'offre (AO) & Procédure formelle par laquelle la faculté sollicite des propositions commerciales auprès de fournisseurs pour l'acquisition de ressources matérielles. \\
\hline
Affectation & Action d'attribuer une ressource matérielle à un département ou à un enseignant précis. \\
\hline
Constat de panne & Document technique rédigé par un technicien décrivant la nature, la fréquence et le type d'une panne. \\
\hline
Fournisseur & Société commerciale externe proposant du matériel informatique en réponse aux appels d'offre de la faculté. \\
\hline
Garantie & Engagement contractuel du fournisseur à réparer ou remplacer le matériel défectueux pendant une durée déterminée à compter de la livraison. \\
\hline
JWT & JSON Web Token. Jeton d'authentification signé numériquement permettant de vérifier l'identité et les droits d'un utilisateur sans session serveur. \\
\hline
Liste noire & Liste des fournisseurs exclus des procédures d'appel d'offre en raison de manquements à leurs engagements. \\
\hline
Moins-disant & Expression désignant le fournisseur ayant soumis l'offre au prix total le plus bas parmi les offres conformes. \\
\hline
Numéro d'inventaire & Identifiant unique attribué à chaque ressource matérielle lors de sa réception, permettant son traçage (également appelé code à barres). \\
\hline
RBAC & Role-Based Access Control. Mécanisme de contrôle d'accès dans lequel les autorisations sont attribuées selon le rôle de l'utilisateur. \\
\hline
Ressource matérielle & Équipement informatique (ordinateur ou imprimante) géré par le système. \\
\hline
Spécifications techniques & Caractéristiques détaillées d'un équipement : pour un ordinateur (CPU, RAM, disque, écran, marque) ; pour une imprimante (vitesse, résolution, marque). \\
\hline
Ticket de panne & Enregistrement dans le système décrivant un dysfonctionnement signalé sur une ressource, de son ouverture à sa clôture. \\
\hline
\end{longtable}

\end{document}
